\PaperPage{39}{5.1 Normalized weights and contraction multipliers}
\label{front:normalized}

Fix a probability measure $\sigma$ on $[0,1]$, a scale $\kappa>0$, and an
anchor $a>0$. In the arithmetic application, these are the actual scale and
the Hausdorff measure supplied under $H_{00}$. The constructions on this
page also make sense for any such measure; their elementary inequalities
do not identify its moments with a prescribed analytic function. Define
\begin{equation}
 w_a(q)=\frac{4\kappa a}{(1+\kappa aq)^2},\qquad
 \lambda_a(q)=q w_a(q)=\frac{4\kappa aq}{(1+\kappa aq)^2}.
\end{equation}
These continuous functions separate the positive density from the variable
which will act by multiplication. This distinction matters: the weight
need not be bounded by one, whereas the multiplier always is.

\begin{proposition}\label{front:kernel-range}
For every $q\in[0,1]$,
\begin{equation}
 0<w_a(q)\leq 4\kappa a,\qquad 0\leq\lambda_a(q)\leq1.
\end{equation}
Moreover, $\lambda_a(q)=1$ if and only if $\kappa aq=1$.
\end{proposition}
\begin{proof}
Set $t=\kappa aq\geq0$. The denominator $(1+t)^2$ is positive and
at least one, proving the weight bounds. The multiplier is nonnegative,
and its upper bound and equality condition follow from the exact identity
\begin{equation}
 1-\lambda_a(q)=\frac{(1-t)^2}{(1+t)^2}.
\end{equation}
Thus equality is possible on the support interval only when
$\kappa a\geq1$, and then only at $q=(\kappa a)^{-1}$.
\end{proof}

For every integer $n\geq0$, multiplication by the nonnegative weight
therefore gives the pointwise domination
\begin{equation}
 0\leq w_a(q)\lambda_a(q)^n\leq w_a(q)\leq4\kappa a.
\end{equation}
This bound is uniform in the order $n$ and in $q$. It needs neither a
density for $\sigma$ nor the exclusion of endpoint atoms. In particular,
the zero endpoint has the explicit values
\begin{equation}
 w_a(0)=4\kappa a,\qquad \lambda_a(0)=0,\qquad
 w_a(0)\lambda_a(0)^n=0\quad(n\geq1).
\end{equation}
An atom at zero contributes to the zeroth coefficient and disappears at
every positive order. Keeping that contribution is necessary when comparing
the base energy with later moments. Conversely, strict positivity of the
weight implies that its integral against a probability measure is strictly
positive. Thus the base vector constructed below cannot vanish, even when
the entire measure is concentrated at an endpoint. If $\kappa a<1$, compactness even gives
$\sup_q\lambda_a(q)<1$; however, the unit bound is the statement valid
uniformly over all positive anchors. The basic kernel inequalities are
audited in \texttt{RouteTwoContractive} and \texttt{ActualPositivePrefix};
the equality case and the following Hilbert interpretation are paper proofs
\cite{source}.

\PaperPage{40}{5.2 A concrete weighted-vector realization}
\label{front:l2}

Work in the complex Hilbert space $\mathcal H=L^2([0,1],\sigma)$, with
inner product linear in its second argument. Put
\begin{equation}
 v_a=\sqrt{w_a},\qquad (T_af)(q)=\lambda_a(q)f(q),\qquad
 \widetilde\mu_n(a)=\int_{[0,1]}w_a(q)\lambda_a(q)^n\,d\sigma(q).
\end{equation}
The tilde distinguishes these measure coefficients from the actual Widder
tower. Since $w_a\leq4\kappa a$ and $\sigma$ has mass one, $v_a$ belongs
to $\mathcal H$. Multiplication respects equality almost everywhere, so
$T_a$ is well defined on the equivalence classes defining $L^2$.

\begin{proposition}\label{front:concrete-orbit}
The operator $T_a$ is a bounded positive self-adjoint contraction, and
\begin{equation}
 \widetilde\mu_n(a)=\langle v_a,T_a^n v_a\rangle,\qquad
 0\leq\widetilde\mu_n(a)\leq\widetilde\mu_0(a)\leq4\kappa a.
\end{equation}
\end{proposition}
\begin{proof}
The range bound on $\lambda_a$ gives
\begin{equation}
 \|T_af\|^2=\int\lambda_a^2|f|^2\,d\sigma\leq\|f\|^2,
 \qquad \langle f,T_af\rangle=\int\lambda_a|f|^2\,d\sigma\geq0.
\end{equation}
Because $\lambda_a$ is real, moving it between the two inner-product
arguments proves self-adjointness. Induction gives $T_a^nf=\lambda_a^nf$.
Substitution of $v_a$ proves the coefficient identity, and integration of
the preceding page's pointwise bounds proves all inequalities.
\end{proof}

The coefficient estimate also follows from a decomposition that survives
in much weaker geometry. For $v=v_a$ and $x=T_a^nv_a$, write
\begin{align}
 \|v\|^4-|\langle v,x\rangle|^2
 &=\underbrace{\|v\|^2(\|v\|^2-\|x\|^2)}_{D_{\rm c}}\nonumber\\
 &\quad+\underbrace{\|v\|^2\|x\|^2-|\langle v,x\rangle|^2}_{D_{\rm g}}.
\end{align}
Contractivity makes $D_{\rm c}$ nonnegative; Cauchy--Schwarz does the same
for $D_{\rm g}$. Thus the defect is the squared Euclidean norm of
$(\sqrt{D_{\rm c}},\sqrt{D_{\rm g}})$. No spectral theorem is needed.

The integral formula supplies more than individual coefficient bounds.
For any polynomial $P(t)=\sum_{j=0}^d c_jt^j$,
\begin{equation}
 \sum_{i,j=0}^d\overline{c_i}c_j\widetilde\mu_{i+j}(a)
 =\int w_a|P(\lambda_a)|^2\,d\sigma\geq0.
\end{equation}
Replacing the weight by $w_a(1-\lambda_a)$ similarly proves positivity
of the matrix of successive moment differences. Both assertions follow
by expanding a finite sum under the integral. They describe additional
structure of this particular realization, beyond the single-coefficient
estimate required at each order by the abstract orbit argument.

The audited \texttt{RealGramGeometry} and \texttt{BKQRContractiveOrbit}
interfaces retain precisely enough scalar inequalities for the real
version of this argument. Their carrier has no required vector-space
structure, and their map $T$ is merely a function with a norm-square
inequality. In particular, those interfaces do not assert linearity,
self-adjointness, completeness, or existence of this $L^2$ model. The
audited defect-square theorem applies to a supplied realization; the
concrete realization above is an additional mathematical construction
\cite{source}.

\PaperPage{41}{5.3 The coefficient resolvent on the unit disk}
\label{front:disk}

The concrete contraction makes a generating function available before any
identification with the arithmetic Widder tower. For a complex parameter
$h$ with $|h|<1$, the operator series converges in norm and gives
\begin{equation}
 (I-hT_a)^{-1}=\sum_{n=0}^{\infty}h^nT_a^n,\qquad
 \|(I-hT_a)^{-1}\|\leq\frac1{1-|h|}.
\end{equation}
Indeed, the operator norm of its tail is bounded by a geometric tail,
and multiplication by $I-hT_a$ telescopes. The same inverse is
multiplication by $(1-h\lambda_a)^{-1}$, since
$|1-h\lambda_a(q)|\geq1-|h|$ everywhere. Consequently
\begin{equation}
 G_a(h):=\langle v_a,(I-hT_a)^{-1}v_a\rangle
 =\int\frac{w_a(q)}{1-h\lambda_a(q)}\,d\sigma(q)
 =\sum_{n=0}^{\infty}\widetilde\mu_n(a)h^n.
\end{equation}

\begin{proposition}\label{front:disk-estimates}
The function $G_a$ is holomorphic on the unit disk. For $0\leq r<1$
and $|h|\leq r$,
\begin{align}
 |G_a(h)|&\leq\frac{\widetilde\mu_0(a)}{1-r},\nonumber\\
 \left|G_a(h)-\sum_{n=0}^{N}\widetilde\mu_n(a)h^n\right|
 &\leq\frac{\widetilde\mu_0(a)r^{N+1}}{1-r}.
\end{align}
For every integer $k\geq0$, its derivatives satisfy
\begin{equation}
 G_a^{(k)}(h)=k!\int
 \frac{w_a\lambda_a^k}{(1-h\lambda_a)^{k+1}}\,d\sigma,
 \qquad |G_a^{(k)}(h)|\leq
 \frac{k!\widetilde\mu_0(a)}{(1-r)^{k+1}}.
\end{equation}
\end{proposition}
\begin{proof}
The coefficient bound from the preceding page supplies a summable
geometric majorant on every closed subdisk. Uniform convergence of
holomorphic partial sums proves holomorphy and the stated remainder
estimate. To differentiate the integral, choose $s$ with $r<s<1$.
Each differentiated integrand on $|h|\leq s$ is bounded by
$k!w_a/(1-s)^{k+1}$, an integrable function. Dominated differentiation
therefore applies. Evaluating the resulting integrand on $|h|\leq r$
gives the sharper displayed bound.
\end{proof}

These estimates are uniform for anchors in any compact positive interval
$[\alpha,\beta]$, because $\widetilde\mu_0(a)\leq4\kappa\beta$.
They establish a stable analytic object with controlled truncation error.
They give no boundary convergence assertion at $|h|=1$, where an atom
with $\lambda_a=1$ contributes a term proportional to $(1-h)^{-1}$.
If $\lambda_a\leq\rho<1$ almost everywhere, the same proof instead
extends the series to $|h|<\rho^{-1}$ when $\rho>0$. When $\rho=0$,
the function is constant. Thus the unit disk is a uniform guarantee,
not a claim that every model has radius exactly one. Coefficient
extraction in that disk gives
\begin{equation}
 G_a^{(n)}(0)=n!\,\widetilde\mu_n(a).
\end{equation}
This recovers the represented moments uniquely from the constructed
analytic function and fixes their normalization without an arbitrary
scalar factor. Finally, $h$ is the coefficient-series
variable; it is not automatically the variable in an actual angular
observable. Such an interpretation requires a separately proved
coefficient identity, to which we now turn.

\PaperPage{42}{5.4 The explicit actual-coefficient obligation}
\label{front:actual-limit}

The archived actual tower is defined from the literal logarithmic
derivative, rather than from a measure. Write $R(a)$ for
\texttt{actualRealLogDerivative} and $\mu_n^{\rm act}(a)$ for
\texttt{actualWidderMoment}. Its definitions are
\begin{equation}
 B(a)=4a\frac{d}{da}\bigl(aR(a)\bigr),\qquad
 \mu_0^{\rm act}=B,\qquad
 \mu_{n+1}^{\rm act}=\mathcal W_n\mu_n^{\rm act},
\end{equation}
where, with $m=n+1$, the specified differential operator is
\begin{equation}
 (\mathcal W_nf)(a)=\frac{2}{m(2m+1)}
 \left(m^2f(a)-af'(a)-a^2f''(a)\right).
\end{equation}
These definitions must be retained when stating a representation theorem.
Replacing their right-hand sides by positive integrals would change the
object requiring identification.

Here $\kappa$ is the actual positive scale. Let $\sigma_N$ denote the finite Bernstein measure of the actual
normalized origin moments, using the repository's grid convention.
Under $H_{00}$ it is a probability measure. The finite quantity called
\texttt{actualKernelApproximation} is
\begin{equation}
 A_{N,n}(a)=\int w_a(q)\lambda_a(q)^n\,d\sigma_N(q).
\end{equation}
The named proposition \texttt{ActualApproximationConvergence} is exactly
the all-order, all-positive-anchor assertion
\begin{equation}\label{front:open-convergence}
 \forall a>0\;\forall n\in\mathbb N,\qquad
 \lim_{N\to\infty}A_{N,n}(a)=\mu_n^{\rm act}(a).
\end{equation}
The source exposes this as a hypothesis; it supplies no inhabitant.
For comparison, weak convergence $\sigma_N\to\sigma$ identifies the
same limit with $\widetilde\mu_n(a)$, because the kernel is continuous
on the compact support interval. Thus identifying the actual differential
tower, rather than merely constructing a measure, is the decisive step.
Uniform bounds on the values of the finite approximants do not, by
themselves, justify passage of first or second derivatives through a
limit. Any argument using the defining recursion must establish the
appropriate regularity and differentiated convergence separately.

\begin{proposition}\label{front:conditional-actual-bound}
Assume $H_{00}$ and \eqref{front:open-convergence}. Then, for $a>0$,
\begin{equation}
 0\leq\mu_n^{\rm act}(a)\leq\mu_0^{\rm act}(a),\qquad
 \sum_{n=0}^{\infty}|\mu_n^{\rm act}(a)|r^n<\infty
 \quad(0\leq r<1).
\end{equation}
\end{proposition}
\begin{proof}
The finite bounds $0\leq A_{N,n}(a)\leq A_{N,0}(a)$ hold at every
depth. Passing both sides to their assumed limits preserves order.
The resulting coefficient bound gives the geometric majorant
$\mu_0^{\rm act}(a)\sum_n r^n$.
\end{proof}

These conditional consequences are audited in \texttt{ActualLimitBridge}
\cite{source}. They do not discharge the original longer angular/Widder
identification program. The completed measure-resolvent route uses exact
origin moments and an analytic defect directly; it reaches its conditional
$H_{00}\!\Rightarrow\!\mathrm{RH}$ endpoint without this additional
coefficient obligation. The distinction prevents a completed route from
being mistaken for a proof of every proposed intermediate construction.

\PaperPage{43}{5.5 Finite data and universal admissibility}
\label{front:finite}

A finite pair of integers can name a mathematical condition that is not
itself a finite computation. The repository's cancellation budget is
precisely such a pair. With $C_n$ the actual row cancellation and $V_n$
the actual variation, its data and validity are
\begin{equation}
 W=(N_0,K)\in\mathbb N^2,\qquad
 \operatorname{Adm}(W):\Longleftrightarrow
 \forall n\geq N_0,\ C_n\leq K V_n.
\end{equation}
The word ``finite'' describes the data fields, whereas admissibility
quantifies over an infinite tail. In particular, nonnegativity of $C_n$
makes the distinguished zero budget equivalent to
\begin{equation}
 H_{00}=\operatorname{Adm}(0,0)
 \quad\Longleftrightarrow\quad\forall n\geq0,\ C_n=0.
\end{equation}
The proved semantic witness theorem is
\begin{equation}
 \bigl(\exists W\in\mathbb N^2,\ \operatorname{Adm}(W)\bigr)
 \quad\Longleftrightarrow\quad\mathrm{RH}.
\end{equation}
Its forward implication uses the cancellation identity and eventual
growth bounds; its reverse implication supplies the zero budget under
RH. Neither direction replaces universal admissibility by evaluation
of a Boolean program \cite{source}.

The separate finite-certificate objective asks for an explicitly defined
type of finite raw data and a terminating checker satisfying
\begin{equation}\label{front:checker-target}
 \bigl(\exists W:\mathrm{FiniteWitness},\
       \mathrm{check}(W)=\mathrm{true}\bigr)
 \quad\Longleftrightarrow\quad\mathrm{RH}.
\end{equation}
Soundness and completeness are logically distinct. A sufficient soundness
theorem would extract an admissible cancellation budget from each
acceptance. Composing this extraction with the semantic theorem would
give RH. Completeness would still require a separate proof that RH
implies the existence of accepted finite data. Neither checker theorem,
nor an unconditional accepted witness, is supplied here.

Finite row computations can establish finitely many inequalities. To
infer an infinite tail, a checker needs an additional proved mechanism,
such as a certified recurrence with an invariant or a verified symbolic
bound whose hypotheses its finite input establishes. This observation is
a specification requirement, not an impossibility claim. A certificate
may encode the finite reason for a universally valid argument, rather
than list every instance of its conclusion. The precise arithmetic
meaning of each encoded expression and the exact actual sequence must
remain part of soundness. Merely placing a proof of the universal
predicate inside a structure would provide mathematical evidence, but
would not by itself define the requested raw-data format and Boolean
acceptance procedure.

\PaperPage{44}{5.6 Conservative formalization targets}
\label{front:targets}

The next formalization tasks can be stated without enlarging the proved
endpoint. First, starting from the probability measure already obtained
under $H_{00}$, construct the complex $L^2$ space, the class of
$\sqrt{w_a}$, and multiplication by $\lambda_a$. The target theorem is
the concrete package
\begin{equation}
 a>0\ \Longrightarrow\
 0\leq T_a=T_a^*,\quad\|T_a\|\leq1,\quad
 \langle v_a,T_a^nv_a\rangle=\widetilde\mu_n(a).
\end{equation}
Here operator positivity means positivity of every quadratic form.
Measurability, square integrability, descent to almost-everywhere
equivalence classes, and boundedness must be proved explicitly. The
resolvent statement then has the conservative form
\begin{equation}
 |h|<1\ \Longrightarrow\
 \langle v_a,(I-hT_a)^{-1}v_a\rangle
 =\int\frac{w_a}{1-h\lambda_a}\,d\sigma.
\end{equation}
Neither theorem mentions the actual Widder tower.

Second, connect this construction to the audited real Gram interface.
One may use real-valued $L^2$ functions for the orbit argument, or take
the real part of the complex inner product on the underlying real
space. The required squared Cauchy inequality then follows from complex
Cauchy--Schwarz. There is also a domain issue: the archived orbit
interface quantifies over every real anchor, while these kernel
inequalities assume $a>0$. A positive-anchor reindexing, or an explicitly
defined zero extension of the model family, must resolve that mismatch.
Such an extension says nothing about actual coefficients at other anchors.

Third, if the longer Widder route is pursued, the additional target is
\begin{equation}
 H_{00}\ \Longrightarrow\
 \forall a>0\;\forall n,\quad
 \mu_n^{\rm act}(a)=\widetilde\mu_n(a).
\end{equation}
A proof must relate the actual logarithmic derivative to the measure
formula and justify the required differentiations and recursion at every
order. Positivity of the model alone cannot replace those identities.

Finally, for a proposed finite checker, define its syntax and evaluation
before seeking the two separate theorem targets
\begin{equation}
 \mathrm{check}(W)=\mathrm{true}\Rightarrow\mathrm{RH},\qquad
 \mathrm{RH}\Rightarrow\exists W,\ \mathrm{check}(W)=\mathrm{true}.
\end{equation}
An accepted certificate would then have exactly the meaning proved by
soundness. These are development targets, not additional audited results.
The paper-level constructions on pages 39--41 do not modify the frozen
Lean sources or their audit boundary at
\texttt{ZETA\_HYPERBOLIC\_ANCHOR\_5}. The distinction between a
mathematical proof presented here and a newly checked Lean declaration
is maintained throughout.
